\section{One life returned to God}
Both readings reject self-sufficiency, but they answer different questions. The first asks what earthly authority and created goods can rightly claim. Jerome's unknowing Cyrus and Chrysostom's limit upon tribute distinguish a genuine service from a claim to divine status. Its practical fruit is ordered allegiance: the freedom to fulfil civic duties without making them absolute.

The second asks how a person becomes capable of such fidelity. Aquinas's inward divine action and Chrysostom's labouring community show why gratitude and effort belong together. Its practical fruit is willing response: the freedom to act without claiming to be the origin of the good one does. Jerome's warning to the gifted teacher brings Cyrus's difficulty into the worshipper's own life.

Either reading becomes incomplete when isolated. An account of obligations can leave the will unchanged; an account of inward renewal can lose sight of real public duties. The coin remains a coin, and the person who renders it remains God's image. The same person needs shelter, a sincere will, purification, nourishment and hope. The Psalm Communion option emphasizes mercy in need; the Mark option emphasizes the saving gift received in faith. The final prayer gives both readings their horizon: present service tends toward a good beyond the power of the servant to produce.
